% !TeX root = ../../../main.tex
\subsubsection{Survival Prediction}\label{ch07:apd:full_survival}

In this section, we report the full concordance analyses for the \textit{sepsis-10} and \textit{sepsis-100} annotation sets in addition to \textit{T2S2} test set, presented in Table~\ref{ch07:tab:concordance_wide}.

{\raggedright
\ChSevenParagraph{Survival analysis: model performance across timepoints and cohorts.}
Table~\ref{ch07:tab:concordance_wide} shows time-dependent concordance results where
instruction-tuned decoder models clearly outperform encoder-based models.
This pattern contrasts with the forecasting tasks (Table~\ref{ch07:tab:forecasting_clean_combined}), where encoders excelled.
Across both the \textit{T2S2} and \textit{sepsis-100} cohorts, decoder models consistently achieve the strongest results.
Notably, \textit{RedPajama-INCITE-7B-Instruct} attains the highest concordance overall (0.76 at 168h, \textit{sepsis-100}) despite lacking any biomedical text in pretraining, while larger models such as \textit{DeepSeek-R1-Llama-70B} and \textit{OLMO-32B-Instruct} also perform robustly.
Encoder models---including biomedically fine-tuned variants like \textit{BioClinical-ModernBERT-base}, which is competitive on \textit{sepsis-100} but not on \textit{T2S2}---are generally surpassed by decoder models.\par}

Predictive performance does not universally improve with more observation time.
Models evaluated on the \textit{T2S2} cohort, often peak at earlier windows (0h or 24h) rather than at 168h.
Meanwhile, the \textit{sepsis-10} cohort experiences clear ceiling effects, with many models achieving perfect concordance (1.00).
These results suggest that, unlike the forecasting tasks, survival prediction in our dataset benefits more from the capabilities inherent in large decoder models than from domain-specific pretrained encoders.

% !TeX root = ../../../main.tex
\begin{table}[htbp]
\caption{Survival analyses evaluated using time-dependent concordance index for each model across \texttt{DeepSeek-R1} test set annotations (\textit{T2S2}, \textit{sepsis-10}, \textit{sepsis-100}), evaluated at three observation time thresholds (0h, 24h, and 168h).
}\ChSevenFloatLabel{ch07:tab:concordance_wide}
\ChSevenTableBodyStart
\begin{ChSevenFloatBody}

\centering
{
\begin{adjustbox}{max width=\linewidth}
\begin{tabular}{lccccccccc}
\toprule
\textbf{Model} 
& \multicolumn{3}{c}{\textit{T2S2}} 
& \multicolumn{3}{c}{\textit{sepsis-10}} 
& \multicolumn{3}{c}{\textit{sepsis-100}} \\
\cmidrule(lr){2-4} \cmidrule(lr){5-7} \cmidrule(lr){8-10}

Time Threshold & 0h & 24h & 168h 
& 0h & 24h & 168h 
& 0h & 24h & 168h \\
\midrule
bert-base-uncased       & 0.60 & 0.61 & 0.53 & 1.00 & 1.00 & 1.00 & 0.58 & 0.62 & 0.58 \\
roberta-base            & 0.55 & 0.60 & 0.60 & 0.75 & 1.00 & 1.00 & 0.55 & 0.58 & 0.57 \\
deberta-v3-small        & 0.56 & 0.57 & 0.57 & 0.75 & 0.00 & 0.75 & 0.52 & 0.49 & 0.51 \\
ModernBERT-base         & 0.52 & 0.58 & 0.58 & 1.00 & 1.00 & 1.00 & 0.54 & 0.51 & 0.51 \\
ModernBERT-large        & 0.53 & 0.59 & 0.58 & 0.50 & 1.00 & 1.00 & 0.59 & 0.55 & 0.49 \\
\textcolor{black}{Bioclinical ModernBERT-base} & 0.57& 0.53& 0.53& 1.00& 0.75& 0.75& 0.63& 0.59& 0.56\\
\textcolor{black}{Bioclinical ModernBERT-large} & 0.54& 0.55 & 0.54 & 0.75 & 0.75 & 0.75& 0.61 & 0.57 & 0.49\\
DeepSeek-R1-Llama-70B   & \textbf{0.64} & \textbf{0.63} & 0.59 & 0.50 & 1.00 & 1.00 & 0.61 & 0.53 & 0.55 \\
Llama-3.3-70B-Instruct  & 0.62 & \textbf{0.63} & 0.58 & 0.50 & 1.00 & 0.75 & 0.64 & 0.57 & 0.53 \\
DeepSeek-R1-Llama-8B    & 0.60 & 0.58 & 0.58 & 1.00 & 1.00 & 1.00 & 0.53 & \textbf{0.66} & {0.65} \\
Llama-3.1-8B-Instruct   & 0.62 & 0.61 & \textbf{0.61} & 1.00 & 1.00 & 1.00 & 0.60 & 0.59 & 0.61 \\
\textcolor{black}{OLMO 32B Instruct} & 0.60& 0.62 & 0.59 & 0.75 & 0.75 & 0.75&\textbf{ 0.66} & 0.58 & 0.56\\
\textcolor{black}{RedPajama-INCITE 7B Instruct} & 0.56&  0.60 & \textbf{0.61} & 0.50 & 1.00 & 1.00& 0.61 & 0.53 & \textbf{0.76}\\
\textcolor{black}{MediPhi-PubMed} & 0.54 & 0.62 & 0.55 & 1.00 & 1.00 & 1.00& 0.62 & 0.58 & 0.61\\
\bottomrule
\end{tabular}
\end{adjustbox}
}



\end{ChSevenFloatBody}
\end{table}
\FloatBarrier
